\begin{pastexam}[label={px:ncu-110-dsa-01}]{中央資工所 110 學年度 複選題第 1 題}
（複選；每題 5 分，全對才給分，答錯一個選項扣 1.5 分）

Which of the following statements are true?
\begin{enumerate}[label=(\Alph*)]
\item When converting the infix expression a*(b+c/d)*(e-f)+g to its postfix form, there are at
most 4 tokens in the stack at any moment.
\item The postfix form of infix expression (a+b-c)*(d-e) is abc-+de-*
\item The infix expression of postfix expression 3k-2m4+*+pab-*/ is ((3-k)+2(m+4))/(p*(a-b))
\item When converting the infix expression a*(b+c/d)*(e-f) to its postfix form, there are at most
3 tokens in the stack at any moment.
\end{enumerate}
\end{pastexam}

\begin{pastsoln}
答案：(A)(C)。

(C) 的中序式寫成 $2(m+4)$，省略了乘號；下面把它讀成 $2*(m+4)$。 (A)、(D) 的「堆疊裡的記號」
包含左括號，這是資結書虛擬碼~6.8 的做法：左括號要推入堆疊，等右括號來時才彈出。

\textbf{(A) 對}。 用資結書第~6.3.3~節的規則轉換：運算元直接輸出；左括號推入；右括號一直彈出到
左括號；運算子先彈出堆疊頂端優先順序\textbf{大於或等於}自己的（遇到左括號停），再推入自己。
\begin{center}
\small
\begin{tabular}{clll}
\toprule
讀到 & 動作 & 堆疊（左邊是底） & 輸出 \\
\midrule
$a$ & 輸出 & & $a$ \\
$*$ & 推入 & $*$ & $a$ \\
$($ & 推入 & $*\ ($ & $a$ \\
$b$ & 輸出 & $*\ ($ & $ab$ \\
$+$ & 頂端是左括號，推入 & $*\ (\ +$ & $ab$ \\
$c$ & 輸出 & $*\ (\ +$ & $abc$ \\
$/$ & $/$ 比 $+$ 高，推入 & $*\ (\ +\ /$ \quad（$4$ 個） & $abc$ \\
$d$ & 輸出 & $*\ (\ +\ /$ & $abcd$ \\
$)$ & 彈出 $/$、$+$，丟掉左括號 & $*$ & $abcd/+$ \\
$*$ & 彈出同級的 $*$，推入 & $*$ & $abcd/+*$ \\
$($ & 推入 & $*\ ($ & $abcd/+*$ \\
$e$ & 輸出 & $*\ ($ & $abcd/+*e$ \\
$-$ & 推入 & $*\ (\ -$ & $abcd/+*e$ \\
$f$ & 輸出 & $*\ (\ -$ & $abcd/+*ef$ \\
$)$ & 彈出 $-$，丟掉左括號 & $*$ & $abcd/+*ef-$ \\
$+$ & 彈出 $*$，推入 & $+$ & $abcd/+*ef-*$ \\
$g$ & 輸出 & $+$ & $abcd/+*ef-*g$ \\
結束 & 彈出剩下的 & & $abcd/+*ef-*g+$ \\
\bottomrule
\end{tabular}
\end{center}
最多的時候是讀到 $/$ 之後的 $4$ 個（$*$、$($、$+$、$/$），所以「最多 $4$ 個」成立。

\textbf{(B) 錯}。 $+$、$-$ 同級、由左往右結合，$a+b-c$ 是 $(a+b)-c$，後序是 $ab+c-$；整個式子是
$ab+c-de-*$。 題目的 $abc-+$ 是 $a+(b-c)$，值雖然一樣，但不是這個中序式的後序形式。

\textbf{(C) 對}。 照資結書虛擬碼~6.7 從左到右掃描，堆疊裡放子運算式：
$3\,k\,-$ 得 $(3-k)$；$2$、$m\,4\,+$ 得 $(m+4)$；$*$ 得 $2*(m+4)$；$+$ 得 $(3-k)+2*(m+4)$；
$p$、$a\,b\,-$ 得 $(a-b)$；$*$ 得 $p*(a-b)$；最後 $/$ 得
\[
\bigl((3-k)+2*(m+4)\bigr)/\bigl(p*(a-b)\bigr),
\]
和題目的式子一樣（只差省略的乘號）。

\textbf{(D) 錯}。 這個式子就是 (A) 去掉最後的 $+g$，前面的步驟完全一樣，讀到 $/$ 時堆疊裡一樣
有 $4$ 個記號，所以「最多 $3$ 個」不成立。
\end{pastsoln}
